\section{JUSTIFICATIVA}
\label{JUSTIFICATIVA}

Até o século XIX, o conhecimento humano evoluía predominantemente de forma cumulativa e gradual. Entretanto, ao longo do século XX, assistiu-se a um avanço acelerado da ciência e da tecnologia, modificando radicalmente os processos produtivos e organizacionais. A incorporação de tecnologias complexas à indústria e a intensificação dos fluxos econômicos globais passaram a exigir profissionais com formação sólida, capazes de acompanhar e promover inovações em contextos altamente dinâmicos.

Nesse cenário, as atividades relacionadas ao controle e à automação de processos assumem papel central no desenvolvimento industrial contemporâneo. A modernização de linhas de produção, a adoção de sistemas inteligentes e a busca por maior produtividade e eficiência operacional têm ampliado a demanda por profissionais especializados, aptos a integrar conhecimentos técnicos avançados com visão sistêmica e capacidade de gestão.

O crescimento industrial verificado no Ceará, especialmente na Região Metropolitana de Fortaleza, reforça essa demanda. A instalação e a expansão de complexos industriais nos setores de siderurgia, metal-mecânica, energia, alimentos, petroquímica e tecnologia intensificaram a necessidade por quadros profissionais altamente qualificados. O Complexo Portuário e Industrial do Pecém e o Polo Industrial de Maracanaú constituem exemplos estratégicos desse movimento, configurando um ambiente produtivo que requer engenheiros com formação atualizada, capacidade de atuação multidisciplinar e domínio dos modernos processos produtivos.


Entretanto, à ocasião elaboração deste PPC, o número de engenheiros formados no Brasil apresentava-se aquém das necessidades do setor produtivo. Segundo o CONFEA, o país formava aproximadamente 40 mil engenheiros por ano, enquanto economias emergentes, como Índia e China, formavam centenas de milhares. Além disso, observava-se elevada taxa de evasão nos cursos de engenharia, relacionada tanto à complexidade das áreas de conhecimento envolvidas quanto às deficiências de base no ensino médio, sobretudo em matemática, física e química. Como consequência, havia déficit anual estimado em aproximadamente 20 mil engenheiros no país, impactando diretamente a capacidade de inovação e o desenvolvimento tecnológico.

{\vm Nos dias atuais..}


A Lei nº 11.892/2008 instituiu a Rede Federal de Educação Profissional, Científica e Tecnológica, atribuindo aos Institutos Federais a responsabilidade de ofertar cursos superiores, incluindo bacharelados em engenharia, como estratégia para o fortalecimento do desenvolvimento econômico e social. Ao longo de sua consolidação, essa Rede passou a distinguir-se pela qualificação de seu corpo docente, pela tradição no ensino técnico e tecnológico, pela inserção territorial em diferentes contextos socioprodutivos e pela articulação orgânica entre ensino, pesquisa e extensão, entendidos como dimensões indissociáveis da formação humana.

Nesse horizonte, a atuação do IFCE tem buscado ampliar o diálogo com os arranjos produtivos locais, reconhecendo que a formação em engenharia deve ultrapassar a mera transmissão de conteúdos técnicos. Tal formação requer o desenvolvimento de capacidades analíticas, pensamento crítico e compreensão do papel social do engenheiro na mediação entre tecnologia, trabalho e sociedade. Assim, pretende-se favorecer a constituição de profissionais capazes de interpretar problemas complexos, propor soluções inovadoras e atuar com responsabilidade ética, social e ambiental.

Nessa perspectiva, o IFCE \textit{Campus} Maracanaú oferta o Curso de Bacharelado em Engenharia de Controle e Automação. O curso orienta-se para a formação de um Bacharel que compreenda os fundamentos científicos e tecnológicos de sua área, seja capaz de acompanhar sua constante evolução e reconheça o impacto de suas decisões no contexto produtivo e social. A formação proposta estimula a prática investigativa, a aprendizagem contínua e o compromisso com o desenvolvimento sustentável, contribuindo para o fortalecimento das dinâmicas locais e regionais e reafirmando a missão institucional do IFCE, em consonância com a Lei nº 11.892/2008.



